% !TEX encoding = UTF-8 Unicode
% ============================================================
%  Trilha Java Básico — Anotações de Aula
%  Aula: Estruturas de Repetição em Java
% ============================================================
\documentclass[12pt,a4paper]{article}

\usepackage[utf8]{inputenc}
\usepackage[T1]{fontenc}
\usepackage[brazil]{babel}
\usepackage{lmodern}
\usepackage{geometry}
\geometry{a4paper, margin=2.5cm}

\usepackage{amsmath,amssymb}
\usepackage{graphicx}
\usepackage{hyperref}
\hypersetup{
  colorlinks=true,
  linkcolor=blue!50!black,
  urlcolor=blue!50!black,
  pdftitle={Aula — Estruturas de Repetição em Java}
}

% --- Caixas para enunciado e solução ---------------------------------
\usepackage[most]{tcolorbox}
\newtcolorbox{enunciado}{
  colback=blue!4, colframe=blue!55!black,
  fonttitle=\bfseries, title=Enunciado do Exemplo,
  arc=2mm, boxrule=0.8pt
}
\newtcolorbox{solucao}{
  colback=green!4, colframe=green!45!black,
  fonttitle=\bfseries, title=Solução,
  arc=2mm, boxrule=0.8pt
}
\newtcolorbox{observacao}{
  colback=yellow!6, colframe=orange!60!black,
  fonttitle=\bfseries, title=Observação,
  arc=2mm, boxrule=0.8pt
}

% --- Código Java ------------------------------------------------------
\usepackage{listings}
\usepackage{xcolor}
\definecolor{codebg}{RGB}{248,248,248}
\definecolor{codecomment}{RGB}{106,153,85}
\definecolor{codestring}{RGB}{163,21,21}
\definecolor{codekeyword}{RGB}{0,0,255}

\lstdefinestyle{java}{
  language=Java,
  backgroundcolor=\color{codebg},
  commentstyle=\color{codecomment}\itshape,
  keywordstyle=\color{codekeyword}\bfseries,
  stringstyle=\color{codestring},
  numberstyle=\tiny\color{gray},
  basicstyle=\ttfamily\footnotesize,
  breaklines=true,
  showstringspaces=false,
  numbers=left,
  numbersep=6pt,
  frame=single,
  rulecolor=\color{gray!40},
  tabsize=4,
  captionpos=b,
  extendedchars=true,
  literate=
    {á}{{\'a}}1 {à}{{\`a}}1 {â}{{\^a}}1 {ã}{{\~a}}1
    {é}{{\'e}}1 {ê}{{\^e}}1 {í}{{\'i}}1 {ó}{{\'o}}1
    {ô}{{\^o}}1 {õ}{{\~o}}1 {ú}{{\'u}}1 {ü}{{\"u}}1
    {ç}{{\c{c}}}1
    {Á}{{\'A}}1 {À}{{\`A}}1 {Â}{{\^A}}1 {Ã}{{\~A}}1
    {É}{{\'E}}1 {Ê}{{\^E}}1 {Í}{{\'I}}1 {Ó}{{\'O}}1
    {Ô}{{\^O}}1 {Õ}{{\~O}}1 {Ú}{{\'U}}1 {Ü}{{\"U}}1
    {Ç}{{\c{C}}}1
}
\lstset{style=java}

% --- Atalhos ----------------------------------------------------------
\newcommand{\codigo}[1]{\lstinline[language=Java,basicstyle=\ttfamily\small]|#1|}

\title{\textbf{Estruturas de Repetição em Java}}
\author{Trilha Java Básico}
\date{Aula de 17 de agosto de 2026}

\begin{document}

\maketitle
\tableofcontents
\newpage

% ============================================================
\section{Introdução}

Estruturas de repetição --- também chamadas de \emph{laços de repetição},
\emph{laços de iteração} ou simplesmente \emph{loops} --- são comandos que
permitem a \textbf{iteração de código}, ou seja, que os comandos presentes em um
bloco sejam repetidos diversas vezes.

Por meio do \emph{controle de fluxo de repetição}, as instruções podem acontecer
recorrentemente (mais de uma vez) até que uma \textbf{condição} determine que
não é mais necessária a continuação da execução daquele bloco.

\begin{observacao}
Os trechos de código abaixo foram reconstruídos a partir da narração da aula.
Valores e nomes exibidos na tela (ex.: elementos de um \emph{array}) podem ter
pequenas variações em relação ao código original do professor, mas a lógica e a
estrutura são fiéis ao que foi explicado.
\end{observacao}

% ============================================================
\section{Classificação das estruturas de repetição}

Existem três estruturas de repetição em Java:

\begin{itemize}
  \item \codigo{for} --- tradução: \textbf{``para''};
  \item \codigo{while} --- tradução: \textbf{``enquanto''};
  \item \codigo{do while} --- tradução: \textbf{``faça enquanto''}.
\end{itemize}

Cada uma tem sua aplicabilidade, suas particularidades estruturais e dicas de
uso, detalhadas a seguir.

% ============================================================
\section{Comando \texttt{for}}

O comando \codigo{for} (``para'') permite que uma \textbf{variável contadora}
seja testada e incrementada a cada iteração. A cada execução do fluxo, essa
variável pode sofrer uma alteração de valor. O \codigo{for} recebe três
informações, definidas na própria chamada do comando:

\begin{enumerate}
  \item a \textbf{variável contadora} (inicialização);
  \item a \textbf{condição} (expressão booleana de validação);
  \item o \textbf{incremento} (atualização da variável).
\end{enumerate}

\subsection{Sintaxe}

\begin{lstlisting}[caption={Estrutura do \texttt{for}}]
for (inicializacao; expressaoBooleana; incremento) {
    // corpo: comandos a serem repetidos
}
\end{lstlisting}

As três partes ficam em uma única linha, separadas por ponto e vírgula
(\codigo{;}):
\begin{itemize}
  \item \textbf{Inicialização} --- declara a variável contadora (ex.: \codigo{int i = 0});
  \item \textbf{Expressão booleana} --- validação; enquanto for \textbf{verdadeira},
        o corpo é executado;
  \item \textbf{Incremento} --- atualiza a variável a cada iteração (ex.: \codigo{i++}).
\end{itemize}

O corpo (bloco entre chaves) é executado inúmeras vezes, enquanto a expressão de
validação for verdadeira. O comando deixa de executar quando a condição se torna
\textbf{falsa}.

O fluxo é: \emph{inicializa} $\rightarrow$ \emph{testa a condição} $\rightarrow$
\emph{executa o corpo} $\rightarrow$ \emph{incrementa} $\rightarrow$
\emph{testa de novo} $\ldots$ até a condição se tornar falsa.

\subsection{Exemplo 1 --- Contador de carneirinhos}

\begin{enunciado}
Joãozinho precisa contar até 20 carneirinhos para pegar no sono. Monte um laço
\codigo{for} com uma variável inteira (\codigo{int}) chamada \codigo{carneirinhos},
que inicia em 1 e, \textbf{enquanto for menor ou igual a 20}, é incrementada de 1
em 1 a cada iteração. A cada repetição, exiba a contagem atual. Ao final, depois
que a condição deixar de ser verdadeira, exiba a mensagem \emph{``Joãozinho
dormiu!''}.
\end{enunciado}

\begin{solucao}
\begin{lstlisting}[caption={ExemploFor.java}]
public class ExemploFor {
    public static void main(String[] args) {
        for (int carneirinhos = 1; carneirinhos <= 20; carneirinhos++) {
            System.out.println("Contando carneirinhos: " + carneirinhos);
        }
        System.out.println("Joãozinho dormiu!");
    }
}
\end{lstlisting}
\end{solucao}

\subsubsection{Simulação passo a passo (depuração)}

Acompanhando a execução (por exemplo, com \emph{breakpoint} na IDE):

\begin{enumerate}
  \item \textbf{Inicialização:} \codigo{carneirinhos = 1}.
  \item \textbf{Teste:} \codigo{1 <= 20} é \emph{verdadeiro} $\rightarrow$ executa o corpo.
  \item \textbf{Corpo:} imprime \emph{``Contando carneirinhos: 1''}.
  \item \textbf{Incremento:} \codigo{carneirinhos++} $\rightarrow$ \codigo{carneirinhos = 2}.
  \item \textbf{Teste:} \codigo{2 <= 20} é \emph{verdadeiro} $\rightarrow$ imprime 2.
  \item $\ldots$ e assim sucessivamente até imprimir 20.
  \item Após imprimir 20, o incremento leva a \codigo{carneirinhos = 21}.
  \item \textbf{Teste:} \codigo{21 <= 20} é \emph{falso} $\rightarrow$ o laço é encerrado.
  \item O programa segue normalmente para \codigo{System.out.println("Joãozinho dormiu!")}.
\end{enumerate}

\begin{observacao}
A saída imprime de 1 a 20 e, ao final, \emph{``Joãozinho dormiu!''}. Isso confirma
que o programa \textbf{continua executando} normalmente após o término do laço.
\end{observacao}

\subsection{Exemplo 2 --- \texttt{for} percorrendo um array}

\begin{enunciado}
Crie um \emph{array} de \codigo{String} com os nomes dos alunos e percorra-o com
um laço \codigo{for}, imprimindo o aluno de cada índice. Use a propriedade
\codigo{length} (tamanho) do \emph{array} como limite da condição.
\end{enunciado}

\begin{solucao}
\begin{lstlisting}[caption={ExemploForArray.java}]
public class ExemploForArray {
    public static void main(String[] args) {
        String[] alunos = { "Felipe", "Jonas", "Julia", "Marcos" };

        for (int x = 0; x < alunos.length; x++) {
            System.out.println("O aluno no índice " + x + " é " + alunos[x]);
        }
    }
}
\end{lstlisting}
\end{solucao}

Explicação linha a linha:
\begin{itemize}
  \item \codigo{String[] alunos = { ... }} --- um \emph{array} é um conjunto de
        elementos do mesmo tipo. Pode ser de \codigo{String}, \codigo{int},
        \codigo{double} ou de qualquer classe (clientes, funcionários etc.).
        Cada elemento fica entre chaves, separado por vírgula.
  \item \codigo{for (int x = 0; x < alunos.length; x++)} --- \codigo{x} representa
        o \textbf{índice} do elemento. O laço continua enquanto \codigo{x} for
        menor que o tamanho do \emph{array}.
  \item \codigo{alunos[x]} --- acessa o aluno que está na posição (índice)
        \codigo{x} da lista.
\end{itemize}

\begin{observacao}
Em \emph{arrays}, o índice dos elementos \textbf{inicia em 0} --- por isso a
variável \codigo{x} começa com 0. Se \codigo{alunos.length} vale 4, os índices
válidos são \codigo{0}, \codigo{1}, \codigo{2} e \codigo{3}.
\end{observacao}

\subsection{\texttt{for each} (``for it'')}

Para interagir com \emph{arrays} e coleções de forma mais agradável, existe o
\codigo{for each}, que percorre diretamente os elementos --- sem a necessidade de
controlar o índice manualmente:

\begin{lstlisting}[caption={Exemplo de \texttt{for each}}]
String[] alunos = { "Felipe", "Jonas", "Julia", "Marcos" };

for (String aluno : alunos) {
    System.out.println("O nome do aluno é: " + aluno);
}
\end{lstlisting}

O uso do \codigo{for each} está fortemente relacionado a cenários com
\emph{arrays}/coleções. A cada iteração, a variável temporária \codigo{aluno}
(do tipo do elemento, aqui \codigo{String}) recebe \textbf{automaticamente} o
elemento atual da coleção. Os dois pontos (\codigo{:}) indicam ``a cada elemento
de''.

% ============================================================
\section{Comandos \texttt{break} e \texttt{continue}}

Dentro de laços, duas palavras têm interação direta com o contexto de repetição:

\begin{itemize}
  \item \codigo{break} --- significa \textbf{quebrar / parar / interromper}.
        Interrompe \textbf{todo} o laço.
  \item \codigo{continue} --- como o nome diz, \textbf{continua o laço}.
        Interrompe somente a \textbf{iteração atual} (pula para a próxima).
\end{itemize}

\begin{enunciado}
Considere um laço \codigo{for} que conta de 1 a 5. Dentro dele existe uma
condição: \emph{se o número for igual a 3, interromper o fluxo}. Qual o resultado
da execução? Em seguida, troque o \codigo{break} por \codigo{continue} e observe
o que muda no resultado.
\end{enunciado}

\begin{solucao}
Com \codigo{break}, imprime apenas \textbf{1 e 2}: ao chegar em 3, o laço é
interrompido por completo --- mesmo a condição determinando a contagem de 1 a 5.

Com \codigo{continue}, imprime \textbf{1, 2, 4 e 5}: ao chegar em 3, apenas
aquela iteração é pulada; o laço segue normalmente para 4 e 5.

\begin{lstlisting}[caption={Com \texttt{break}}]
for (int numero = 1; numero <= 5; numero++) {
    if (numero == 3)
        break;
    System.out.println(numero);
}
// Saída: 1, 2
\end{lstlisting}

\begin{lstlisting}[caption={Com \texttt{continue}}]
for (int numero = 1; numero <= 5; numero++) {
    if (numero == 3)
        continue;
    System.out.println(numero);
}
// Saída: 1, 2, 4, 5
\end{lstlisting}
\end{solucao}

\begin{observacao}
O \codigo{break} é útil quando uma \textbf{regra de negócio}, ao se tornar
verdadeira, deve interromper toda a execução. O \codigo{continue} é útil quando se
quer apenas \emph{desconsiderar} uma determinada iteração, sem parar o laço.
\end{observacao}

% ============================================================
\section{Comando \texttt{while}}

O laço \codigo{while} (``enquanto'') determina que, \textbf{enquanto uma condição
for válida, o bloco de código será executado}. Ele \textbf{testa a condição antes
de executar o código}: se a condição já for inválida no primeiro teste, o bloco
\textbf{nem chega a ser executado}.

\subsection{Sintaxe}

\begin{lstlisting}[caption={Estrutura do \texttt{while}}]
while (expressaoBooleana) {
    // corpo: comandos executados enquanto a condição for verdadeira
}
\end{lstlisting}

O comando será executado até que a expressão de validação se torne \textbf{falsa}.

\begin{observacao}
É preciso garantir que, em algum momento, a condição se torne \textbf{falsa}.
Caso contrário, cria-se um \textbf{laço infinito} (o bloco nunca para de
executar).
\end{observacao}

\subsection{Exemplo 3 --- A mesada do Joãozinho}

\begin{enunciado}
Joãozinho recebeu \textbf{R\$ 50,00} de mesada e foi a uma loja de doces para
gastar todo o dinheiro. Enquanto o valor da mesada for \textbf{maior que zero},
ele adiciona ao carrinho um doce de valor aleatório (entre R\$ 2 e R\$ 8) e
desconta esse valor da mesada. Ao final, exiba o valor restante da mesada e uma
mensagem informando que Joãozinho gastou toda a sua mesada em doces.
\end{enunciado}

\begin{solucao}
\begin{lstlisting}[caption={ExemploWhile.java}]
import java.util.concurrent.ThreadLocalRandom;

public class ExemploWhile {
    public static void main(String[] args) {
        double mesada = 50.0;

        while (mesada > 0) {
            double valorDoce = valorAleatorio();
            if (valorDoce > mesada) {
                valorDoce = mesada;   // adaptação: evita mesada negativa
            }
            System.out.println("Doce do valor: " + valorDoce + " adicionado no carrinho");
            mesada = mesada - valorDoce;
        }

        System.out.println("Mesada: " + mesada);
        System.out.println("Joãozinho gastou toda a sua mesada em doces");
    }

    private static double valorAleatorio() {
        return ThreadLocalRandom.current().nextDouble(2, 8);
    }
}
\end{lstlisting}
\end{solucao}

Explicação:
\begin{itemize}
  \item \codigo{double mesada = 50.0;} --- a mesada inicia com 50.
  \item \codigo{while (mesada > 0)} --- enquanto a mesada for maior que zero, o
        laço executa: pega um doce de valor aleatório e subtrai da mesada.
  \item \codigo{valorAleatorio()} --- método auxiliar que retorna um \codigo{double}
        aleatório entre 2 e 8, usando \codigo{ThreadLocalRandom}. (Não é requisito
        do \codigo{while}; serve apenas para simular valores aleatórios.)
  \item \codigo{if (valorDoce > mesada) valorDoce = mesada;} --- adaptação para
        que a mesada não fique negativa (regra de negócio). Sem esse trecho o
        código também funciona, mas poderia terminar com valor negativo.
\end{itemize}

\begin{observacao}
A quantidade de iterações não é fixa: depende dos valores aleatórios sorteados.
O laço executa enquanto a condição for verdadeira --- neste exemplo, até a mesada
chegar a zero (ou próximo disso) e a condição \codigo{mesada > 0} se tornar falsa.
\end{observacao}

% ============================================================
\section{Comando \texttt{do while}}

O \codigo{do while} (``faça enquanto'') tem a mesma proposta do \codigo{while}
(executa enquanto a condição for válida), mas com uma diferença fundamental:
\textbf{testa a condição depois de executar o código}, garantindo que o bloco seja
executado \textbf{pelo menos uma vez}.

Assim, mesmo que a condição já seja inválida no primeiro teste, o bloco executa
uma vez antes de parar.

\subsection{Sintaxe}

\begin{lstlisting}[caption={Estrutura do \texttt{do while}}]
do {
    // corpo: executado ao menos uma vez
} while (expressaoBooleana);
\end{lstlisting}

Note o ponto e vírgula (\codigo{;}) ao final da linha do \codigo{while} ---
obrigatório nessa estrutura.

\subsection{Exemplo 4 --- O telefone tocando}

\begin{enunciado}
Joãozinho ligou para um amigo. Enquanto o telefone estiver \textbf{tocando}, ele
aguarda; quando o amigo \textbf{atender}, o fluxo é interrompido e parte-se para a
conversa (\emph{``Alô!''}). Use um \codigo{do while} para simular a chamada: o
telefone deve ``tocar'' ao menos uma vez antes de verificar se o amigo atendeu.
\end{enunciado}

\begin{solucao}
\begin{lstlisting}[caption={ExemploDoWhile.java}]
import java.util.Random;

public class ExemploDoWhile {
    public static void main(String[] args) {
        System.out.println("Discando...");

        do {
            System.out.println("Telefone tocando");
        } while (tocando());

        System.out.println("Alô!!!");
    }

    private static boolean tocando() {
        boolean atendeu = new Random().nextInt(3) + 1 == 1;   // valor aleatório entre 1 e 3
        System.out.println("Atendeu? " + atendeu);
        return !atendeu;   // continua "tocando" enquanto NÃO atendeu
    }
}
\end{lstlisting}
\end{solucao}

Explicação:
\begin{itemize}
  \item \codigo{do { ... } while (tocando());} --- o corpo executa ao menos
        uma vez (``Telefone tocando'') e, só depois, a condição é testada.
  \item \codigo{tocando()} --- método auxiliar que retorna \codigo{boolean}
        (\emph{sim ou não}: atendeu ou não atendeu). Gera um valor aleatório entre
        1 e 3; se for igual a 1, significa que \textbf{atendeu}.
  \item \codigo{return !atendeu;} --- usa o operador unário de negação
        (\codigo{!}, ``exclamação''). Se atendeu, nega para \codigo{false} e o
        laço para de ``tocar''; senão, retorna \codigo{true} e o telefone continua
        tocando.
\end{itemize}

\begin{observacao}
Por ser aleatório, o número de ``toques'' muda a cada execução: às vezes o amigo
atende de primeira; outras vezes o telefone toca várias vezes. O que é garantido é
que o telefone \textbf{toca ao menos uma vez} antes da verificação --- o que não
aconteceria com o \codigo{while}, que testa a condição antes.
\end{observacao}

% ============================================================
\section{Comparação: \texttt{while} \textit{vs} \texttt{do while}}

\begin{center}
\begin{tabular}{p{0.46\textwidth} p{0.46\textwidth}}
  \hline
  \textbf{\codigo{while}} & \textbf{\codigo{do while}} \\
  \hline
  Testa a condição \textbf{antes} de executar. & Testa a condição \textbf{depois}
  de executar. \\
  Pode \textbf{não executar nenhuma vez} (se a condição já for falsa no primeiro
  teste). & Executa \textbf{pelo menos uma vez}, mesmo que a condição seja falsa
  de início. \\
  \hline
\end{tabular}
\end{center}

% ============================================================
\section{Resumo e dicas de prática}

\begin{itemize}
  \item \codigo{for} --- ideal quando se conhece (ou se controla) a quantidade de
        repetições por meio de um contador.
  \item \codigo{while} --- ideal quando a repetição depende de uma condição que é
        validada antes de cada execução.
  \item \codigo{do while} --- ideal quando o bloco \textbf{precisa} executar ao
        menos uma vez antes da validação.
  \item \codigo{break} interrompe todo o laço; \codigo{continue} pula apenas a
        iteração atual.
  \item Em \emph{arrays}, o índice começa em \codigo{0}.
\end{itemize}

Para consolidar o conteúdo, o professor recomenda praticar de forma isolada, por
exemplo:
\begin{itemize}
  \item um contador de 0 a 100;
  \item um contador de números positivos;
  \item definir cenários onde \codigo{break} e \codigo{continue} se encaixam;
  \item criar um \emph{array} e listar todas as frutas de uma cesta de frutas;
  \item usar a depuração (\emph{breakpoint}) para acompanhar o valor das variáveis
        a cada iteração.
\end{itemize}

\end{document}
